\section{Related work}
\label{sec:related}

\paragraph{What an explanation should satisfy}
Interpretability lacks a single definition~\cite{lipton2018,doshivelez2017}.
Jacovi and Goldberg~\cite{jacovi2020} separate faithfulness, whether an
explanation reflects the model's reasoning, from plausibility, whether it
convinces a person, and warn that evaluations often measure the second. Chen et
al.~\cite{chen2020} draw the related line between explanations true to the
model and true to the data, and Janzing et al.~\cite{janzing2020} argue that
feature relevance is a causal question about interventions on the model's
inputs. Our on/off contributions are interventions of exactly that kind, made
on the plan inputs the merchant controls.

\paragraph{Feature attributions and their failures}
LIME~\cite{ribeiro2016} and SHAP~\cite{lundberg2017}, with its exact form for
trees~\cite{lundberg2020}, are the most used attributions~\cite{adadi2018}.
Saliency maps can barely change when a network's weights are
randomised~\cite{adebayo2018}, move with input shifts the model
ignores~\cite{kindermans2019}, differ for near-identical
inputs~\cite{alvarezmelis2018}, and rank features whose removal does not hurt
the model~\cite{hooker2019}. Rudin~\cite{rudin2019} argues for models that are
interpretable by construction. The world model and the lift model are of that
kind for the plan, which lets this audit compare a built-in explanation with
post-hoc ones on the same forecasts.

\paragraph{Explanations of time series and forecasts}
Rojat et al.~\cite{rojat2021} survey explanation methods for time series.
Ismail et al.~\cite{ismail2020} and Tonekaboni et
al.~\cite{tonekaboni2020} show on generated data with known important time steps
that saliency methods designed for images often fail on temporal data, the
approach our known-answer study takes for retail plans. Among forecasting
models, the Temporal Fusion Transformer~\cite{tft} reports attention weights
as explanations. We know of no study that checks the explanation of a demand
forecast against a measured promotional or price effect.

\paragraph{Promotion and price effects}
Marketing science has measured promotional effects for decades: the split of a
promotion's sales into brand switching, timing and quantity~\cite{gupta1988,
vanheerde2003}, the dip after a promotion~\cite{mace2004,vanheerde2000}, and
the size of price elasticities and what drives it~\cite{bijmolt2005}. Promotion forecasting systems built on these
effects include PromoCast~\cite{promocast}; Ma, Fildes and Huang~\cite{ma2016}
and Fildes et al.~\cite{fildes2022} review the forecasting side. These elasticities describe temporary cuts. Regular price changes move far fewer
units: in Dominick's randomised store-level pricing experiments~\cite{hoch1994},
regular prices were raised or lowered 10\% across many categories, and volume
moved about 3\% over the following 16 weeks, an elasticity near $-0.3$. Hoch et al.~\cite{hoch1995} and
Montgomery~\cite{montgomery1997} model store-level price response on the same
data.

\paragraph{Quasi-experiments in scanner data}
The public Dominick's files do not mark the experimental stores or weeks.
Abrams, Gui and Horta\c{c}su~\cite{abrams} reconstructed them with the
original documentation and found that some published price changes did not
replicate. We therefore treat the price deviations visible in the public data
as a quasi-experiment, never as the experiments, and estimate their effect by
difference in differences~\cite{angrist2009} with standard errors clustered by
store~\cite{bertrand2004,cameron2015}.
